\section{Completed Results}
\label{sec:completed-results}

\subsection{Canonical implementation validation}
Table~\ref{tab:canonical-engine-benchmark} summarizes the independent pricing check.
All 36 controlled cases pass at both node counts, with no integration warnings,
no-arbitrage failures, or parity failures. Agreement supports implementation validity
on the controlled domain; it does not prove universal correctness or equivalence to
unavailable teammate code.

\begin{table}[htbp]
\centering
\caption{Independent adaptive-quadrature validation of the production Double Heston pricer.}
\label{tab:canonical-engine-benchmark}
\begin{tabular}{lrr}
\toprule
Diagnostic & 64-node production & 96-node production \\
\midrule
Controlled cases passing & 36/36 & 36/36 \\
RMSE versus reference & $5.458 \times 10^{-13}$ & $4.223 \times 10^{-12}$ \\
MAE versus reference & $5.184 \times 10^{-13}$ & $4.064 \times 10^{-12}$ \\
Maximum absolute difference & $8.100 \times 10^{-13}$ & $5.699 \times 10^{-12}$ \\
No-arbitrage failures & 0 & 0 \\
Parity maximum error & $7.11 \times 10^{-15}$ & $7.11 \times 10^{-15}$ \\
\bottomrule
\end{tabular}
\par\vspace{0.35em}
\begin{minipage}{0.94\linewidth}\small The reference independently implements the affine characteristic function and uses adaptive SciPy quadrature; it does not import the production pricer.\end{minipage}
\end{table}


\subsection{Development-market BS/Heston/DH comparison}
Table~\ref{tab:real-market-development} reports the NTPC development pilot. Heston has
a slightly better holdout price RMSE than Double Heston, while Double Heston has a
slightly better holdout IV RMSE. Under the predeclared 5\%-improvement rule the result
is \texttt{NO\_CLEAR\_WINNER}. Both Heston and Double Heston selected capped iterates;
the DH multi-start classification is materially unstable. This is development context,
not final market performance.

\begin{table}[htbp]
\centering
\caption{Development-market BS/Heston/DH calibration pilot for NTPC on 2026-07-15.}
\label{tab:real-market-development}
\begin{tabular}{lcccccc}
\toprule
Model & Parameters & Calibration price RMSE & Holdout price RMSE & Calibration IV RMSE & Holdout IV RMSE & Runtime (s) \\
\midrule
BLACK\_SCHOLES & 1 & 0.326902 & 1.05334 & 0.0118538 & 0.0804675 & 0.0118729 \\
HESTON & 5 & 0.238578 & 0.910569 & 0.00790466 & 0.0729962 & 41.2127 \\
DOUBLE\_HESTON & 10 & 0.234321 & 0.926825 & 0.00778971 & 0.0725916 & 228.607 \\
\bottomrule
\end{tabular}
\par\vspace{0.35em}
\begin{minipage}{0.94\linewidth}\small This is a development diagnostic, not a final real-market result. The declared winner rule returned NO\_CLEAR\_WINNER; fitted DH parameters are not claimed to be true market parameters.\end{minipage}
\end{table}


\subsection{Frozen synthetic benchmark}
Table~\ref{tab:primary-comparison} presents seed-mean neural metrics and the frozen
traditional representatives. Traditional calibration attains the best parameter RMSE
and essentially exact repricing, but averages about 376 seconds per surface in the
recorded environment. Both networks evaluate orders of magnitude faster and emit valid
vectors on all surfaces, but have range-scaled parameter RMSE near 0.166.
Figure~\ref{fig:speed-accuracy} visualizes this tradeoff without treating hardware
runtime as quality.

\begin{table}[htbp]
\centering
\caption{Frozen synthetic-test comparison on 1{,}250 untouched known-truth surfaces.}
\label{tab:primary-comparison}
\begin{tabular}{lccc}
\toprule
Metric & Model 1 ANN & Model 2 informed & Traditional \\
\midrule
Range-scaled parameter RMSE & 0.165896 & 0.166443 & 0.115122 \\
Standardized parameter RMSE & 0.901305 & 0.904044 & 0.600244 \\
v0 total MAE & 0.010059 & 0.010367 & 0.00056 \\
theta total MAE & 0.033022 & 0.033017 & 0.014561 \\
Constraint validity & 100.00\% & 100.00\% & 100.00\% \\
Mean repricing nRMSE & 0.000727 & 0.000789 & $5.699 \times 10^{-9}$ \\
Repricing p95 nRMSE & 0.001968 & 0.002106 & $2.365 \times 10^{-8}$ \\
Repricing $\leq 10^{-4}$ & 2.00\% & 0.96\% & 100.00\% \\
Repricing $\leq 10^{-3}$ & 77.84\% & 78.08\% & 100.00\% \\
Parameter RMSE $\leq 0.25$ & 91.20\% & 91.04\% & 94.00\% \\
Mean evaluation seconds per surface & $3.708 \times 10^{-6}$ & $8.142 \times 10^{-6}$ & 375.96 \\
\bottomrule
\end{tabular}
\par\vspace{0.35em}
\begin{minipage}{0.94\linewidth}\small Neural values are means over seeds 11, 22, and 33. Traditional uses its frozen representative start. Inference and calibration runtimes are hardware-specific provenance and are not a quality score.\end{minipage}
\end{table}


\begin{figure}[htbp]
\centering
\includegraphics[width=0.78\linewidth]{speed_accuracy_tradeoff.pdf}
\caption{Speed/accuracy position of each method on the untouched synthetic test split.
Runtime axes are hardware-specific provenance.}
\label{fig:speed-accuracy}
\end{figure}

\subsection{Pricing fit does not certify parameters}
Traditional calibration reprices within $10^{-4}$ on 100\% of test surfaces, yet only
$71.92\%$ also have range-scaled parameter RMSE $\leq0.10$ and $94.00\%$ have RMSE
$\leq0.25$. Among neural predictions with repricing nRMSE $\leq10^{-3}$, only about
$7.6\%$ additionally satisfy parameter RMSE $\leq0.10$, although roughly $95\%$ remain
within the looser 0.25 threshold. Figure~\ref{fig:fit-not-recovery} shows why the two
metrics must remain separate.

\begin{figure}[htbp]
\centering
\includegraphics[width=0.70\linewidth]{fit_is_not_recovery.pdf}
\caption{Mean repricing error versus known-truth parameter error. Distance along the
repricing axis does not determine distance along the recovery axis.}
\label{fig:fit-not-recovery}
\end{figure}

\subsection{Stability}
Neural headline metrics are stable across seeds, but per-surface prediction dispersion
is larger for Model2 than Model1; adding differentiable repricing did not reduce seed
dispersion in the completed primary evidence. Traditional multi-start behavior is much
more dispersed: the per-surface start-disagreement rate is $89.16\%$, and 3{,}192 of
3{,}750 starts reach the frozen function-evaluation cap.

\subsection{Observation-noise status}
Table~\ref{tab:completed-neural-noise} and Figure~\ref{fig:neural-noise} report only
the completed neural full-population evaluations. They do not imply that the three-way
positive-noise comparison is finished. Positive-level traditional subset calibration
remains active/pending, so no three-way noise conclusion is drawn here.

\begin{table}[htbp]
\centering
\caption{Completed neural full-test-population observation-noise evaluations.}
\label{tab:completed-neural-noise}
\begin{tabular}{lccccc}
\toprule
Noise & M1 parameter RMSE & M1 clean-latent nRMSE & M2 parameter RMSE & M2 clean-latent nRMSE & Mean recovery $\leq0.25$ \\
\midrule
0\% & 0.165991 & 0.000789 & 0.166914 & 0.001027 & 91.00\% \\
0.10\% & 0.165992 & 0.000789 & 0.166915 & 0.001027 & 91.00\% \\
0.25\% & 0.16599 & 0.000789 & 0.166913 & 0.001027 & 90.99\% \\
0.50\% & 0.166 & 0.000791 & 0.166923 & 0.001029 & 91.00\% \\
1.00\% & 0.166028 & 0.000797 & 0.166948 & 0.001033 & 90.99\% \\
\bottomrule
\end{tabular}
\par\vspace{0.35em}
\begin{minipage}{0.94\linewidth}\small Seed-mean values over seeds 11, 22, and 33 for all 1,250 test surfaces at each level, read from pinned commit e58b36d1ff6c. Positive-level traditional subset completion remains pending and is excluded here.\end{minipage}
\end{table}


\begin{figure}[htbp]
\centering
\includegraphics[width=0.76\linewidth]{completed_neural_noise.pdf}
\caption{Completed seed-mean neural recovery rates across observation-noise levels.
Traditional positive-noise results are omitted because they are incomplete.}
\label{fig:neural-noise}
\end{figure}
